\documentclass[12pt,a4paper]{report}
\usepackage[utf8]{inputenc}
\usepackage[T1]{fontenc}
\usepackage[french]{babel}
\usepackage{lmodern}
\usepackage[margin=2.5cm]{geometry}
\usepackage{amsmath,amssymb}
\usepackage{booktabs}
\usepackage{array}
\usepackage{enumitem}
\usepackage{xcolor}
\usepackage{colortbl}
\usepackage{listings}
\usepackage{tikz}
\usetikzlibrary{positioning}
\usepackage{pgfplots}
\pgfplotsset{compat=1.16}
\usepackage{caption}
\usepackage{float}
\usepackage{etoolbox}
\usepackage[colorlinks=true,linkcolor=blue!70!black,urlcolor=blue!70!black]{hyperref}

\hypersetup{pdftitle={Benchmark ABR et tas}, pdfauthor={BENNOUR RAFAA}}
\captionsetup{labelsep=endash, font=small}
\captionsetup[table]{position=above}
\emergencystretch=2em
\widowpenalty=10000 \clubpenalty=10000
\setlist{itemsep=2pt}

% ---------------- Code ----------------
\definecolor{fondcode}{HTML}{F6F6F3}
\definecolor{bordcode}{HTML}{CCCCCC}
\definecolor{motcle}{HTML}{1F5FA8}
\definecolor{commentaire}{HTML}{6E6E6E}
\definecolor{ok}{HTML}{D4EDD4}
\definecolor{ko}{HTML}{F6D1CC}
\definecolor{cmain}{HTML}{1F5FA8}
\definecolor{cpq}{HTML}{D35400}
\lstset{
  language=C++, basicstyle=\ttfamily\footnotesize,
  keywordstyle=\color{motcle}\bfseries, commentstyle=\color{commentaire}\itshape,
  backgroundcolor=\color{fondcode}, frame=single, rulecolor=\color{bordcode},
  columns=fullflexible, keepspaces=true, showstringspaces=false,
  xleftmargin=4pt, xrightmargin=4pt, aboveskip=8pt, belowskip=8pt,
  literate={é}{{\'e}}1 {è}{{\`e}}1 {à}{{\`a}}1 {ê}{{\^e}}1 {ç}{{\c{c}}}1
           {ù}{{\`u}}1 {É}{{\'E}}1 {î}{{\^i}}1 {ô}{{\^o}}1 {û}{{\^u}}1 {â}{{\^a}}1
}
\lstdefinestyle{sortie}{language={}, basicstyle=\ttfamily\scriptsize, keywordstyle={}, backgroundcolor=\color{white}}
\BeforeBeginEnvironment{lstlisting}{\par\noindent\begin{minipage}{\linewidth}}
\AfterEndEnvironment{lstlisting}{\end{minipage}\par}

% ---------------- Arbres ----------------
\tikzset{
  noeud/.style={circle, draw=black!70, fill=blue!8, minimum size=7mm, inner sep=0pt, font=\small},
  racine/.style={noeud, fill=orange!30},
  level distance=11mm,
}

% =====================================================================
\begin{document}

\begin{titlepage}
  \centering
  \vspace*{2cm}
  {\huge Benchmark Comparatif\par}
  \vspace{0.4cm}
  {\LARGE Arbre binaire de recherche et tas\par}
  \vspace{0.3cm}
  {\Large Tas max écrit à la main vs \texttt{priority\_queue}\par}
  \vspace{1.3cm}
  {\large Insertion, recherche, suppression et extraction de la racine :\\[2pt]
  analyse de performance comparative\par}
  \vspace{1.3cm}
  \begin{tikzpicture}[level 1/.style={sibling distance=34mm}, level 2/.style={sibling distance=17mm}]
    \node[racine] {67}
      child { node[noeud] {57}
        child { node[noeud] {28} }
        child { node[noeud] {40} } }
      child { node[noeud] {10}
        child { node[noeud] {8} }
        child { node[noeud] {9} } };
  \end{tikzpicture}
  \vspace{1.3cm}

  \begin{tabular}{r@{\quad}l}
    \textbf{Auteur :} & BENNOUR RAFAA \\
    \textbf{Formation :} & ILISI -- 2\textsuperscript{e} année \\
    \textbf{Langage :} & C++17 \\
    \textbf{Compilateur :} & g++ 15.2 \texttt{-O2} \\
  \end{tabular}
  \vfill
\end{titlepage}

\setcounter{tocdepth}{1}
\tableofcontents

% =====================================================================
\chapter{Introduction}

\section{Contexte et objectifs}
Ce rapport étudie deux structures en forme d'\textbf{arbre binaire} (chaque nœud a au plus deux enfants), qui rangent les données de deux façons différentes :
\begin{description}[leftmargin=0.5cm, style=nextline]
  \item[Arbre binaire de recherche (ABR)]
  pour chaque nœud, les valeurs plus petites sont à gauche et les plus grandes à droite. On peut donc \textbf{chercher} une valeur en descendant dans l'arbre, comme dans un dictionnaire.
  \item[Tas (heap)]
  chaque parent est plus grand (\textbf{tas max}) ou plus petit (\textbf{tas min}) que ses enfants. On ne cherche pas une valeur quelconque : on veut seulement obtenir très vite \textbf{le maximum} (ou le minimum), qui est toujours à la racine.
\end{description}

La partie principale du travail demandé porte sur le \textbf{tas max}, réalisé de deux manières :
\begin{enumerate}
  \item \textbf{à la main, sans conteneur} : un simple tableau \texttt{int*} alloué avec \texttt{new}, et nos propres fonctions \texttt{monter} et \texttt{descendre} ;
  \item \textbf{avec \texttt{std::priority\_queue<int>}}, l'adaptateur de la STL, qui est un tas max par défaut.
\end{enumerate}
Les deux versions sont comparées sur les deux opérations d'un tas : \textbf{l'insertion} et \textbf{l'extraction de la racine} (la plus grande valeur).

\subsection{Objectifs de l'étude}
\begin{enumerate}
  \item Implémenter un ABR complet : insertion, recherche, minimum, maximum, suppression (3 cas) et les quatre parcours.
  \item Montrer l'effet de \textbf{l'ordre d'insertion} sur la hauteur de l'ABR, et donc sur sa vitesse.
  \item Implémenter un tas max et un tas min sans conteneur.
  \item Comparer le tas max écrit à la main avec \texttt{priority\_queue} : qui insère le plus vite, et qui extrait la racine le plus vite ?
\end{enumerate}

\section{Les deux structures en image}
Les mêmes valeurs ne donnent pas du tout le même arbre :

\begin{figure}[H]
  \centering
  \begin{minipage}{0.48\linewidth}\centering
  \begin{tikzpicture}[level 1/.style={sibling distance=30mm}, level 2/.style={sibling distance=15mm}, level 3/.style={sibling distance=9mm}]
    \node[racine] {50}
      child { node[noeud] {30}
        child { node[noeud] {20} }
        child { node[noeud] {40}
          child { node[noeud] {35} }
          child { node[noeud] {45} } } }
      child { node[noeud] {70}
        child { node[noeud] {60}
          child[missing] {}
          child { node[noeud] {65} } }
        child { node[noeud] {80} } };
  \end{tikzpicture}\\[2pt]
  \small ABR : gauche $<$ nœud $<$ droite
  \end{minipage}\hfill
  \begin{minipage}{0.48\linewidth}\centering
  \begin{tikzpicture}[level 1/.style={sibling distance=30mm}, level 2/.style={sibling distance=15mm}, level 3/.style={sibling distance=8mm}]
    \node[racine] {67}
      child { node[noeud] {57}
        child { node[noeud] {28} child { node[noeud] {1} } child { node[noeud] {2} } }
        child { node[noeud] {40} child { node[noeud] {7} } child { node[noeud] {29} } } }
      child { node[noeud] {10}
        child { node[noeud] {8} child { node[noeud] {6} } child[missing] {} }
        child { node[noeud] {9} } };
  \end{tikzpicture}\\[2pt]
  \small Tas max : parent $\geq$ enfants
  \end{minipage}
  \caption{À gauche l'ABR de la démonstration, à droite le tas max construit avec 7 9 67 1 28 8 10 29 2 57 40 6.}
\end{figure}

Un tas est toujours un arbre \textbf{complet} (rempli niveau par niveau, de gauche à droite). C'est pour cela qu'on peut le ranger dans un simple tableau, sans pointeurs :
\begin{center}
\begin{tikzpicture}[x=8.5mm, y=8mm]
  \foreach \v [count=\i from 0] in {67,57,10,28,40,8,9,1,2,7,29,6} {
    \draw[fill=blue!8] (\i,0) rectangle ++(1,1);
    \node at (\i+0.5,0.5) {\small\v};
    \node[font=\scriptsize, text=black!60] at (\i+0.5,-0.3) {\i};
  }
  \draw[fill=orange!30] (0,0) rectangle ++(1,1); \node at (0.5,0.5) {\small 67};
\end{tikzpicture}\\[4pt]
\small Pour la case $i$ : parent $=(i-1)/2$, enfant gauche $=2i+1$, enfant droit $=2i+2$.
\end{center}

\section{Jeux de données}
\begin{itemize}
  \item \textbf{Démonstrations} : ABR avec 50 30 70 20 40 60 80 35 45 65 ; tas avec 7 9 67 1 28 8 10 29 2 57 40 6.
  \item \textbf{Benchmark ABR} : les clés $0,1,\dots,N-1$, insérées soit dans l'ordre \textbf{croissant}, soit dans un ordre \textbf{aléatoire} (mélange avec une graine fixe). $N$ va jusqu'à $20\,000$ en ordre croissant (au-delà, c'est trop lent) et jusqu'à $1\,000\,000$ en ordre aléatoire.
  \item \textbf{Benchmark tas} : $N = 10^3$ à $10^7$ entiers, dans trois ordres : \textbf{aléatoire}, \textbf{croissant} et \textbf{décroissant}.
\end{itemize}

\section{Opérations mesurées}
\begin{center}\small
\begin{tabular}{lll}
  \toprule
  Structure & Opération & Ce qu'on mesure \\
  \midrule
  ABR & \texttt{inserer} & insérer les $N$ clés \\
      & \texttt{rechercher} & chercher chacune des $N$ clés \\
      & \texttt{supprimer} & supprimer les $N$ clés une par une \\
  \midrule
  Tas max & insertion & \texttt{inserer(x)} / \texttt{pq.push(x)}, $N$ fois \\
          & extraction de la racine & \texttt{extraireRacine()} / \texttt{top()} + \texttt{pop()} \\
  \bottomrule
\end{tabular}
\end{center}
Chaque résultat est un \textbf{temps moyen par opération}, en nanosecondes (ns).

\section{Complexités théoriques}
\begin{table}[H]
\centering
\caption{Complexités ($h$ = hauteur de l'ABR, $n$ = nombre d'éléments).}
\small
\begin{tabular}{lccc}
  \toprule
  Opération & ABR équilibré & ABR dégénéré (liste) & Tas \\
  \midrule
  Insertion & $O(\log n)$ & $O(n)$ & $O(\log n)$ \\
  Recherche d'une valeur & $O(\log n)$ & $O(n)$ & $O(n)$ \\
  Lire le maximum & $O(\log n)$ & $O(n)$ & $O(1)$ (tas max) \\
  Supprimer / extraire le maximum & $O(\log n)$ & $O(n)$ & $O(\log n)$ \\
  \bottomrule
\end{tabular}
\end{table}
Pour l'ABR, tout dépend de la \textbf{hauteur} $h$ : chaque opération coûte $O(h)$. Si les valeurs arrivent dans un ordre aléatoire, $h$ reste proche de $\log_2 n$ ; si elles arrivent triées, chaque nouvelle valeur va à droite de la précédente et l'arbre devient une \textbf{liste} : $h = n$. Le tas, lui, est toujours complet : sa hauteur vaut toujours $\lfloor\log_2 n\rfloor$, quel que soit l'ordre.

% =====================================================================
\chapter{Méthodologie}

\section{Environnement}
\begin{itemize}
  \item Processeur : Intel Core i7-8550U à 1,8~GHz (4 cœurs, 8 threads), cache L2 de 256~Ko par cœur, cache L3 de 8~Mo.
  \item Mémoire : 8~Go de DDR4 à 2400~MHz.
  \item Système : Windows 11 Pro ; compilateur g++ 15.2 (MinGW-w64), options \texttt{-O2 -std=c++17}.
  \item C'est un ordinateur portable : sa fréquence varie (mode Turbo, température). Pour le vérifier, le benchmark du tas a été lancé deux fois. Les temps changent de 10 à 30~\% d'une série à l'autre, mais les tendances restent les mêmes. Les tableaux montrent la première série.
\end{itemize}

\section{Protocole de mesure}
\begin{itemize}
  \item Chronomètre : \texttt{std::chrono::steady\_clock}.
  \item Chaque mesure est répétée (5 fois pour le tas, 3 fois pour l'ABR et pour $N=10^7$) et l'on garde \textbf{le meilleur temps}, le moins perturbé par le système.
  \item Les deux versions du tas reçoivent \textbf{exactement les mêmes données}.
  \item Les valeurs extraites sont additionnées dans une variable \texttt{volatile} pour que le compilateur ne supprime pas le calcul.
  \item \textbf{Vérification} avant les mesures : sur $100\,000$ valeurs, le tas écrit à la main et \texttt{priority\_queue} doivent sortir les mêmes valeurs, dans l'ordre décroissant. Pour l'ABR, chaque recherche doit réussir et chaque suppression doit trouver sa clé.
\end{itemize}

\section{Implémentation}

\subsection{Arbre binaire de recherche}
Un nœud contient une valeur et deux pointeurs. L'insertion et la recherche sont \textbf{itératives} (une boucle, pas de récursion) : en ordre croissant l'arbre a une hauteur de $20\,000$, et une fonction récursive risquerait de saturer la pile.
\begin{lstlisting}
struct Noeud { int valeur; Noeud* gauche = nullptr; Noeud* droite = nullptr; };

bool rechercher(int v) const {
    Noeud* c = racine;
    while (c != nullptr) {
        if (v == c->valeur) return true;
        c = (v < c->valeur) ? c->gauche : c->droite;   // gauche si plus petit
    }
    return false;
}
\end{lstlisting}

La \textbf{suppression} traite trois cas :
\begin{enumerate}
  \item le nœud est une \textbf{feuille} : on l'enlève ;
  \item il a \textbf{un seul enfant} : l'enfant prend sa place ;
  \item il a \textbf{deux enfants} : on copie dans le nœud le \textbf{minimum du sous-arbre droit} (la valeur juste au-dessus), puis on supprime ce minimum, qui n'a jamais d'enfant gauche.
\end{enumerate}
\begin{lstlisting}
Noeud* n = *p;                       // p : adresse du lien qui pointe vers le noeud
if (n->gauche == nullptr)      { *p = n->droite; delete n; }   // cas 1 et 2
else if (n->droite == nullptr) { *p = n->gauche; delete n; }   // cas 2
else {                                                         // cas 3
    Noeud** q = &n->droite;
    while ((*q)->gauche != nullptr) q = &(*q)->gauche;         // minimum à droite
    Noeud* m = *q;
    n->valeur = m->valeur;
    *q = m->droite;
    delete m;
}
\end{lstlisting}
Les parcours préfixe, infixe et suffixe sont récursifs ; le parcours en largeur utilise une file (\texttt{queue}).

\subsection{Tas max sans conteneur}
Le tas est un tableau \texttt{int* t} géré à la main : quand il est plein, on alloue un tableau deux fois plus grand et on recopie (\texttt{agrandir}).
\begin{lstlisting}
void inserer(int v) {                  // ajouter à la fin, puis faire monter
    if (n == capacite) agrandir();
    t[n] = v;
    int i = n++;
    while (i > 0 && t[i] > t[(i-1)/2]) { swap(t[i], t[(i-1)/2]); i = (i-1)/2; }
}

int extraireRacine() {                 // le dernier remplace la racine, puis descend
    int max = t[0];
    t[0] = t[--n];
    descendre(0);                      // échange avec le plus grand enfant
    return max;
}
\end{lstlisting}

\subsection{Tas min}
Le tas min est le même code avec les comparaisons inversées (\texttt{<} au lieu de \texttt{>}) : chaque parent est plus petit que ses enfants, et la racine est le minimum.

\subsection{Tas max avec \texttt{priority\_queue}}
\begin{lstlisting}
#include <queue>
priority_queue<int> pq;                // tas max par défaut (un vector en dessous)
for (int x : donnees) pq.push(x);                     // insertion
while (!pq.empty()) { s += pq.top(); pq.pop(); }      // extraction de la racine
\end{lstlisting}

\subsection{Organisation du projet}
\begin{center}
\begin{tabular}{ll}
  \toprule
  Fichier & Rôle \\
  \midrule
  \texttt{ArbreBinaireRecherche.h} & classe \texttt{ABR} \\
  \texttt{TasMax.h}, \texttt{TasMin.h} & tas écrits à la main (tableau \texttt{new}/\texttt{delete}) \\
  \texttt{DemoABR.cpp}, \texttt{DemoTas.cpp} & démonstrations sur de petits exemples \\
  \texttt{BenchmarkABR.cpp} & ordre croissant vs aléatoire \\
  \texttt{BenchmarkTas.cpp} & tas à la main vs \texttt{priority\_queue} \\
  \bottomrule
\end{tabular}
\end{center}

% =====================================================================
\chapter{Résultats : arbre binaire de recherche}

\section{Démonstration}
\begin{lstlisting}[style=sortie]
Nombre de noeuds : 10      Hauteur : 4      Minimum : 20      Maximum : 80
Parcours prefixe (racine, gauche, droite) : 50 30 20 40 35 45 70 60 65 80
Parcours infixe  (gauche, racine, droite) : 20 30 35 40 45 50 60 65 70 80
Parcours suffixe (gauche, droite, racine) : 20 35 45 40 30 65 60 80 70 50
Parcours en largeur (niveau par niveau)   : 50 30 70 20 40 60 80 35 45 65
Recherche de 45 : trouve      Recherche de 55 : absent
Insertion de 40 (deja present) : refuse
Suppression de 20 (feuille)           -> 30 35 40 45 50 60 65 70 80
Suppression de 60 (un seul enfant 65) -> 30 35 40 45 50 65 70 80
Suppression de 50, la racine (deux enfants, remplacee par 65) -> 30 35 40 45 65 70 80
\end{lstlisting}
Le parcours \textbf{infixe donne toujours les valeurs triées} : c'est la propriété qui définit un ABR. Après chaque suppression, l'infixe reste trié, ce qui prouve que l'arbre est resté correct.

\section{Données brutes}
\begin{table}[H]
\centering
\caption{ABR, insertion dans l'ordre \textbf{croissant} (ns par opération).}
\begin{tabular}{rrrrr}
  \toprule
  $N$ & Hauteur & Insertion & Recherche & Suppression \\
  \midrule
  1\,000  & 1\,000  & 1\,866  & 2\,045  & 920 \\
  2\,000  & 2\,000  & 4\,441  & 5\,752  & 2\,223 \\
  5\,000  & 5\,000  & 11\,475 & 14\,627 & 6\,572 \\
  10\,000 & 10\,000 & 23\,921 & 29\,734 & 15\,482 \\
  20\,000 & 20\,000 & 48\,991 & 65\,438 & 31\,689 \\
  \bottomrule
\end{tabular}
\end{table}

\begin{table}[H]
\centering
\caption{ABR, insertion dans un ordre \textbf{aléatoire} (ns par opération).}
\begin{tabular}{rrrrrr}
  \toprule
  $N$ & Hauteur & $\lceil\log_2 N\rceil$ & Insertion & Recherche & Suppression \\
  \midrule
  1\,000     & 22 & 10 & 144.2  & 38.1  & 104.5 \\
  2\,000     & 23 & 11 & 131.6  & 43.4  & 111.3 \\
  5\,000     & 29 & 13 & 169.8  & 62.0  & 135.6 \\
  10\,000    & 29 & 14 & 265.8  & 93.0  & 190.2 \\
  20\,000    & 33 & 15 & 338.9  & 150.2 & 274.2 \\
  50\,000    & 38 & 16 & 352.2  & 163.7 & 288.3 \\
  100\,000   & 41 & 17 & 462.2  & 180.0 & 331.2 \\
  1\,000\,000 & 50 & 20 & 1\,049.9 & 720.0 & 979.8 \\
  \bottomrule
\end{tabular}
\end{table}

\section{Visualisation}
\begin{figure}[H]
\centering
\begin{tikzpicture}
\begin{loglogaxis}[width=0.85\linewidth, height=7cm, xlabel={$N$}, ylabel={ns par insertion},
  legend pos=south east, grid=major, grid style={black!12}]
  \addplot[cpq, thick, mark=square*] coordinates {(1000,1865.9) (2000,4440.9) (5000,11474.8) (10000,23920.7) (20000,48990.5)};
  \addplot[cmain, thick, mark=*] coordinates {(1000,144.2) (2000,131.6) (5000,169.8) (10000,265.8) (20000,338.9) (50000,352.2) (100000,462.2) (1000000,1049.9)};
  \legend{ordre croissant, ordre aléatoire}
\end{loglogaxis}
\end{tikzpicture}
\caption{Temps d'insertion dans l'ABR (échelles logarithmiques).}
\end{figure}

\section{Analyse}
\begin{itemize}
  \item \textbf{Ordre croissant} : la hauteur est égale à $N$, l'arbre est une liste. Quand $N$ double, le temps par opération double aussi (un peu plus, de 2 à 2,4 fois) : chaque opération est en $O(n)$, et remplir l'arbre coûte $O(n^2)$. À $N=20\,000$, l'insertion est \textbf{145 fois} plus lente qu'en ordre aléatoire (48\,991 contre 339~ns).
  \item \textbf{Ordre aléatoire} : la hauteur reste petite (50 pour un million de nœuds, environ $2{,}5\times\log_2 N$). Le temps reste très petit par rapport à l'ordre croissant, mais il augmente \textbf{plus vite que la hauteur} : de $10^3$ à $10^6$, la hauteur est multipliée par 2,3 et le temps d'insertion par 7.
  \item La raison est le cache. Un nœud occupe environ 32 octets : à $N=10^5$ l'arbre fait environ 3~Mo et tient dans le cache L3 (8~Mo) ; à $N=10^6$ il fait environ 32~Mo et n'y tient plus. Le temps d'insertion passe alors de 462 à 1\,050~ns alors que la hauteur n'augmente que de 41 à 50 : les nœuds sont dispersés en mémoire, donc chaque descente vers un enfant est un accès lent à la mémoire.
\end{itemize}

% =====================================================================
\chapter{Résultats : tas max à la main vs \texttt{priority\_queue}}

\section{Démonstration : tas max et tas min}
\begin{lstlisting}[style=sortie]
===== TAS MAX =====
Tableau : 67 57 10 28 40 8 9 1 2 7 29 6
Racine (maximum) : 67
Extractions successives : 67 57 40 29 28 10 9 8 7 6 2 1

===== TAS MIN =====
Tableau : 1 2 6 7 28 8 10 29 9 57 40 67
Racine (minimum) : 1
Extractions successives : 1 2 6 7 8 9 10 28 29 40 57 67
\end{lstlisting}
Extraire la racine plusieurs fois donne les valeurs \textbf{triées} : dans l'ordre décroissant pour le tas max, croissant pour le tas min (c'est le principe du tri par tas). Le tableau, lui, n'est pas trié : seule la règle parent/enfants est respectée.

\section{Données brutes}
\begin{table}[H]
\centering
\caption{Tas max : temps moyen par opération (ns). En vert, la version la plus rapide.}
\newcommand{\g}{\cellcolor{ok}}
\begin{tabular}{llrrrr}
  \toprule
  & & \multicolumn{2}{c}{Insertion} & \multicolumn{2}{c}{Extraction de la racine} \\
  \cmidrule(lr){3-4}\cmidrule(lr){5-6}
  Données & $N$ & À la main & \texttt{priority\_queue} & À la main & \texttt{priority\_queue} \\
  \midrule
  aléatoire & $10^3$ & \g 20.1 & 26.4 & 86.7  & \g 50.1 \\
            & $10^4$ & \g 25.9 & 28.1 & 103.2 & \g 55.4 \\
            & $10^5$ & \g 24.7 & 26.7 & 147.3 & \g 81.1 \\
            & $10^6$ & 23.9 & \g 22.3 & 229.6 & \g 187.4 \\
            & $10^7$ & 28.7 & \g 26.8 & \g 399.9 & 624.4 \\
  \midrule
  croissant & $10^3$ & \g 16.4 & 19.4 & 64.3  & \g 29.8 \\
            & $10^4$ & \g 26.3 & 31.1 & 74.4  & \g 47.6 \\
            & $10^5$ & \g 44.5 & 50.1 & 107.6 & \g 79.4 \\
            & $10^6$ & \g 44.0 & 55.8 & 106.1 & \g 90.2 \\
            & $10^7$ & \g 50.1 & 64.7 & 126.5 & \g 107.6 \\
  \midrule
  décroissant & $10^3$ & \g 6.6 & 8.2  & 81.4  & \g 49.7 \\
              & $10^4$ & 15.0 & \g 14.7 & 108.4 & \g 66.4 \\
              & $10^5$ & 11.4 & \g 9.3  & 106.8 & \g 66.2 \\
              & $10^6$ & 7.0  & \g 6.0  & 120.0 & \g 99.6 \\
              & $10^7$ & 9.2  & \g 7.3  & 145.1 & \g 121.6 \\
  \bottomrule
\end{tabular}
\end{table}

\section{Visualisation}
\begin{figure}[H]
\centering
\begin{tikzpicture}
\begin{semilogxaxis}[width=0.85\linewidth, height=6.2cm, xlabel={$N$}, ylabel={ns par insertion},
  legend pos=north west, grid=major, grid style={black!12}, ymin=0]
  \addplot[cmain, thick, mark=*] coordinates {(1e3,20.1) (1e4,25.9) (1e5,24.7) (1e6,23.9) (1e7,28.7)};
  \addplot[cpq, thick, mark=square*] coordinates {(1e3,26.4) (1e4,28.1) (1e5,26.7) (1e6,22.3) (1e7,26.8)};
  \legend{à la main, \texttt{priority\_queue}}
\end{semilogxaxis}
\end{tikzpicture}
\caption{Insertion, données aléatoires : la version à la main est devant jusqu'à $10^5$, puis les deux sont presque égales.}
\end{figure}

\begin{figure}[H]
\centering
\begin{tikzpicture}
\begin{semilogxaxis}[width=0.85\linewidth, height=6.2cm, xlabel={$N$}, ylabel={ns par extraction},
  legend pos=north west, grid=major, grid style={black!12}, ymin=0]
  \addplot[cmain, thick, mark=*] coordinates {(1e3,86.7) (1e4,103.2) (1e5,147.3) (1e6,229.6) (1e7,399.9)};
  \addplot[cpq, thick, mark=square*] coordinates {(1e3,50.1) (1e4,55.4) (1e5,81.1) (1e6,187.4) (1e7,624.4)};
  \legend{à la main, \texttt{priority\_queue}}
\end{semilogxaxis}
\end{tikzpicture}
\caption{Extraction de la racine, données aléatoires : \texttt{priority\_queue} est devant jusqu'à $10^6$, les courbes se croisent entre $10^6$ et $10^7$.}
\end{figure}

\section{Analyse}

\subsection{Insertion : la version à la main gagne le plus souvent}
Le tas écrit à la main est plus rapide dans \textbf{9 cas sur 15} :
\begin{itemize}
  \item en ordre \textbf{croissant}, il gagne pour toutes les tailles, de 11 à 23~\% ;
  \item en ordre \textbf{aléatoire}, il gagne jusqu'à $10^5$ (de 7 à 24~\%). À $10^6$ et $10^7$, \texttt{priority\_queue} est un peu devant (environ 7~\%), mais ce n'est pas stable : dans la deuxième série, le tas à la main gagne à $10^6$ et les deux sont à égalité à $10^7$ ;
  \item en ordre \textbf{décroissant}, les temps sont très petits (6 à 15~ns) et les écarts sont du même ordre que les variations d'une série à l'autre : on ne peut pas désigner de gagnant.
\end{itemize}
Les deux font le même travail (ajouter à la fin puis faire monter), mais notre version est plus simple :
\begin{itemize}
  \item \texttt{priority\_queue::push} passe par \texttt{vector::push\_back} puis par l'algorithme général \texttt{push\_heap}, écrit pour n'importe quel type et n'importe quel comparateur ;
  \item notre \texttt{inserer} travaille directement sur un tableau d'\texttt{int}, avec une comparaison \texttt{>} écrite en dur.
\end{itemize}
L'ordre des données compte beaucoup plus que la version choisie :
\begin{itemize}
  \item \textbf{décroissant} (6 à 15~ns) : chaque nouvelle valeur est plus petite que son parent, elle ne monte pas du tout. C'est le \textbf{meilleur cas} : $O(1)$ par insertion.
  \item \textbf{croissant} : chaque nouvelle valeur est la plus grande et monte \textbf{jusqu'à la racine}. C'est le \textbf{pire cas} : $\log_2 n$ échanges à chaque fois. On le voit bien à partir de $10^5$ (44 à 50~ns contre 24 à 29~ns en aléatoire). À $N=10^3$, l'ordre croissant est pourtant un peu plus rapide que l'aléatoire (16 contre 20~ns) : probablement parce que la comparaison est toujours vraie, donc le processeur la prévoit parfaitement, alors qu'en aléatoire il se trompe souvent.
  \item \textbf{aléatoire} : en moyenne, une valeur ne monte que d'un ou deux niveaux. L'insertion coûte donc presque $O(1)$ en moyenne, ce qui explique les temps presque constants.
\end{itemize}

\subsection{Extraction de la racine : \texttt{priority\_queue} gagne presque partout}
\begin{itemize}
  \item \texttt{priority\_queue} est plus rapide dans \textbf{14 cas sur 15}. Pour les petits tas l'écart est grand : 50,1 contre 86,7~ns à $N=10^3$ en aléatoire, soit 42~\% de moins, et 29,8 contre 64,3~ns en croissant.
  \item En aléatoire, l'écart reste autour de 45~\% jusqu'à $10^5$, puis il diminue nettement : à $10^6$, 187 contre 230~ns (18~\%).
  \item Seule exception : à $N=10^7$ en aléatoire, le tas à la main est devant (400 contre 624~ns). Ce résultat se répète dans la deuxième série (468 contre 812~ns), ce n'est donc pas une erreur de mesure isolée.
\end{itemize}

\textbf{Pourquoi \texttt{priority\_queue} est-elle plus rapide ?} \texttt{pop\_heap} de la bibliothèque standard (libstdc++) n'utilise pas la même méthode que notre \texttt{descendre} :
\begin{center}
\begin{tabular}{p{0.44\linewidth}p{0.44\linewidth}}
  \toprule
  Notre \texttt{descendre} & \texttt{pop\_heap} de la STL \\
  \midrule
  On place le dernier élément à la racine, puis on l'\textbf{échange} avec le plus grand enfant, niveau après niveau. &
  On crée un \textbf{trou} à la racine et on y fait remonter le plus grand enfant, jusqu'en bas ; puis on place le dernier élément dans le trou et on le fait remonter un peu. \\
  2 comparaisons par niveau + 1 échange (3 écritures) & 1 comparaison par niveau + 1 simple copie \\
  S'arrête dès que l'élément est bien placé & Descend toujours jusqu'en bas \\
  \bottomrule
\end{tabular}
\end{center}
L'élément venu de la fin du tableau est petit : il redescend presque toujours tout en bas. Le test « est-ce que je m'arrête ? » de notre version est donc presque toujours inutile, et la méthode du trou fait moins de travail par niveau.

Sur cette machine, avec g++ 15.2, ce travail en moins suffit à rendre \texttt{priority\_queue} plus rapide même pour les petits tas, qui tiennent entièrement dans le cache.

\textbf{Et l'exception à $10^7$ ?} À $N=10^7$, le tas occupe 40~Mo, cinq fois plus que le cache L3 (8~Mo, soit environ deux millions d'\texttt{int}). Mais la taille seule n'explique pas l'exception : en ordre croissant et décroissant, le tas fait aussi 40~Mo à $N=10^7$, et \texttt{priority\_queue} reste devant.

Ce qui change avec des données aléatoires, c'est que le chemin suivi à chaque extraction est imprévisible : dans le bas de l'arbre, chaque niveau coûte un accès lent à la mémoire et une erreur de prédiction du processeur. Le temps dépend alors surtout de ces accès, et les deux méthodes en font à peu près autant : l'avantage de la méthode du trou disparaît. Cela explique que l'écart se réduise dès $10^6$, mais pas pourquoi \texttt{priority\_queue} devient \textbf{nettement plus lente} à $10^7$. Nous n'avons pas d'explication vérifiée pour ce dernier point : il faudrait un outil de profilage pour compter les accès mémoire de chaque version.

% =====================================================================
\chapter{Analyse comparative approfondie}

\section{ABR ou tas : chacun son usage}
\begin{table}[H]
\centering
\caption{Comparaison des deux structures.}
\begin{tabular}{p{0.32\linewidth}p{0.29\linewidth}p{0.29\linewidth}}
  \toprule
  & ABR & Tas \\
  \midrule
  Règle & gauche $<$ nœud $<$ droite & parent $\geq$ enfants (tas max) \\
  Forme & dépend de l'ordre d'insertion & toujours complète \\
  Hauteur & entre $\log_2 n$ et $n$ & toujours $\lfloor\log_2 n\rfloor$ \\
  Chercher une valeur & oui, $O(h)$ & non (il faut tout parcourir) \\
  Obtenir le maximum & descendre à droite, $O(h)$ & la racine, $O(1)$ \\
  Valeurs triées & parcours infixe & extractions successives \\
  Stockage & nœuds + pointeurs & un tableau \\
  Usage typique & dictionnaire, ensemble trié & file de priorité, tri par tas, Dijkstra \\
  \bottomrule
\end{tabular}
\end{table}
L'ordre d'insertion est le point faible de l'ABR simple (145 fois plus lent en ordre croissant) ; le tas n'a pas ce problème, car il reste toujours complet. Pour avoir un arbre de recherche toujours rapide, il faut un arbre \textbf{équilibré} (AVL, arbre rouge-noir) : c'est ce qu'utilise \texttt{std::map}.

\section{Mémoire}
\begin{itemize}
  \item \textbf{ABR} : chaque nœud contient un \texttt{int} (4 octets) et deux pointeurs (2 × 8 octets), soit 24 octets avec l'alignement, plus le surcoût de \texttt{new} : environ \textbf{32 octets par valeur}. Les nœuds sont dispersés en mémoire.
  \item \textbf{Tas} : \textbf{4 octets par valeur}, rangées les unes à côté des autres. Le tableau double quand il est plein, donc au pire la moitié des cases est vide.
  \item Le tas à la main et \texttt{priority\_queue} utilisent \textbf{la même mémoire} : notre tableau double comme le \texttt{vector} interne de \texttt{priority\_queue}.
\end{itemize}
Un tas utilise donc environ 8 fois moins de mémoire qu'un ABR pour les mêmes valeurs, et il profite beaucoup mieux du cache.

\section{Tas à la main ou \texttt{priority\_queue} ?}
\begin{table}[H]
\centering
\caption{Bilan de la comparaison.}
\begin{tabular}{p{0.34\linewidth}p{0.28\linewidth}p{0.28\linewidth}}
  \toprule
  Critère & À la main & \texttt{priority\_queue} \\
  \midrule
  Insertion & \cellcolor{ok} plus rapide le plus souvent (9 cas sur 15) & \\
  Extraction de la racine & & \cellcolor{ok} plus rapide (14 cas sur 15) \\
  Extraction, $N=10^7$ aléatoire & \cellcolor{ok} plus rapide & \\
  Quantité de code & une classe entière & \cellcolor{ok} une ligne \\
  Risque d'erreur & \texttt{new}/\texttt{delete}, indices & \cellcolor{ok} code déjà testé \\
  Changer d'ordre (tas min) & réécrire les comparaisons & \cellcolor{ok} \texttt{greater<int>} \\
  Comprendre le fonctionnement & \cellcolor{ok} tout est visible & caché \\
  \bottomrule
\end{tabular}
\end{table}

% =====================================================================
\chapter{Conclusion}

\section{Synthèse des résultats}
\begin{enumerate}
  \item L'ABR fonctionne correctement : le parcours infixe reste trié après les trois cas de suppression.
  \item Sa vitesse dépend de sa hauteur : avec des valeurs triées, il devient une liste et il est \textbf{145 fois plus lent} à $N=20\,000$.
  \item Le tas reste toujours complet, donc toujours rapide ; il donne le maximum (ou le minimum) en $O(1)$.
  \item Le tas max écrit à la main \textbf{insère plus vite} que \texttt{priority\_queue} dans 9 cas sur 15 (de 7 à 24~\%), surtout avec des données croissantes ; pour les données décroissantes, les écarts sont trop faibles pour conclure.
  \item Pour \textbf{l'extraction de la racine}, \texttt{priority\_queue} est plus rapide dans 14 cas sur 15 (jusqu'à 42~\% de moins), grâce à la méthode du trou, qui fait moins de travail par niveau. Seule exception : $N=10^7$ en données aléatoires, où le tas à la main est devant ; nous n'en avons pas d'explication vérifiée.
\end{enumerate}

\section{Recommandations}
\begin{itemize}
  \item Pour \textbf{chercher} des valeurs : un ABR, mais équilibré (\texttt{std::set}, \texttt{std::map}) si l'ordre des données n'est pas aléatoire.
  \item Pour obtenir sans cesse \textbf{le plus grand} (ou le plus petit) : un tas.
  \item Dans un vrai programme : \texttt{priority\_queue}, plus courte, plus sûre, et plus rapide pour l'extraction.
  \item Écrire le tas à la main reste utile pour \textbf{apprendre} comment il fonctionne ; son seul vrai avantage mesuré ici est une insertion un peu plus rapide.
\end{itemize}

\section{Conclusion finale}
L'ABR et le tas sont deux arbres binaires, mais ils ne répondent pas à la même question : l'ABR permet de \textbf{chercher} n'importe quelle valeur, le tas donne seulement \textbf{la plus grande}, mais toujours en temps constant. Pour le tas max, le code écrit à la main insère un peu plus vite, mais la technique plus fine de la bibliothèque standard l'emporte nettement pour l'extraction. Les mesures dépendent aussi de la machine et du compilateur : sur un autre PC, avec g++ 13, le tas à la main extrayait plus vite jusqu'à un million d'éléments. Aucune des deux versions n'est la plus rapide partout : le résultat dépend de l'opération, de la taille, de l'ordre des données, de la machine et du compilateur. C'est pour cela qu'il faut mesurer plutôt que deviner.

\end{document}
